\section{Tiny ImageNet: Architecture and Evaluation Protocol}
\label{app:modern_protocol}

\paragraph{Data and frozen representations.}
We use the standard Tiny ImageNet-200 training and labelled validation splits:
100,000 training images (500 per class) and 10,000 validation images (50 per
class), available from the
\href{http://cs231n.stanford.edu/tiny-imagenet-200.zip}{standard source distribution}.
The frozen backbone is
torchvision's ImageNet-pretrained ConvNeXt-Tiny~\citep{liu2022convnext}, with
the preprocessing supplied with its default pretrained weights; no
additional training augmentation is applied. The residual head receives
768-dimensional pooled features. The attention head receives the $7\times7$
spatial feature map after the frozen backbone's final channel LayerNorm,
stored as 49 tokens of dimension 768. Backbone weights are never updated.
The torchvision implementation is distributed under the BSD 3-Clause licence;
the Tiny ImageNet and ImageNet source terms apply to the images and pretrained
weights. We do not redistribute these assets in the manuscript source package.

\paragraph{Residual SwiGLU head.}
A learned projection maps pooled features to $h_0\in\R^{512}$. The tied
recurrence is
\[
 h_{t+1}=h_t+\frac{\lambda}{T_{\rm train}}
 D\!\left(W_o\left[\operatorname{SiLU}(W_g\operatorname{LN}(h_t)+b_g)
 \odot(W_v\operatorname{LN}(h_t)+b_v)\right]+b_o\right),
\]
where the gate and value projections have width 2048 and $D$ is dropout
with probability 0.1 during training. A shared LayerNorm and linear
200-class head classify the resulting state. There is no input reinjection
after $h_0$ and no intermediate supervision.

\paragraph{Shared-attention head.}
A learned projection maps each spatial token to width 256. A learned class
token and fixed two-dimensional sinusoidal spatial positions are added once
before recursion. Each tied block applies pre-normalised eight-head
self-attention, followed by a pre-normalised SwiGLU feed-forward sublayer
of width 1024. Each residual branch is scaled by $\lambda/T_{\rm train}$;
attention probabilities, attention outputs, and feed-forward outputs use
dropout 0.1. Attention is unmasked over the 50 tokens. A shared LayerNorm
and linear classifier read the class token. Parameters are shared across
all recurrent steps, without intermediate supervision or input reinjection.

\paragraph{Training and checkpoint selection.}
Each architecture is trained separately at $T_{\rm train}=8$ and 32, with
paired seeds 101, 102, and 103: 12 completed training runs. All runs use
terminal cross-entropy, AdamW (learning rate $3\times10^{-4}$, weight decay
0.01), 30 epochs, two warm-up epochs followed by cosine decay, label smoothing
0.1, gradient clipping at norm 1, and nominal batch size 512. Attention uses
physical batches of 128 and four accumulation steps, retaining the final
partial accumulation group. Runs use BF16 on an RTX 5080 Laptop GPU.
Checkpoints are selected by terminal validation accuracy at the training
horizon and gain 1, with the earliest maximum retained. The same validation
set supplies reported evaluation accuracies; there is no additional untouched
test split. These are controlled representation experiments, not benchmark
leaderboard estimates.

\paragraph{Horizon comparison.}
All final-panel evaluations fix gain $\lambda=1$ and use horizons 8, 16, 24,
and 32. Within a trained model, the step scale remains fixed as the evaluation
horizon changes. Between the eight- and 32-step training conditions, both
training depth and the prescribed residual scale $1/T_{\rm train}$ change.
The matched condition therefore tests the complete horizon-dependent training
configuration; it does not isolate training depth from step size. Early-eight
features average the first eight post-update states (class-token states for
attention) and are decoded by the shared trained head. Terminal loss is
$D=A(8)-A(32)$ in percentage points. The paired interaction is
$D_{T_{\rm train}=8}-D_{T_{\rm train}=32}$ for each seed. All $\pm$ values in
these additional experiments denote sample standard deviation across the
three seeds, with denominator $n-1$. No significance test is used at $n=3$.

\begin{table}[htbp]
\centering
\caption{\textbf{Final four-condition panel.} Accuracy is in percent; loss is $A(8)-A(32)$ in percentage points. Means and sample standard deviations across three seeds.}
\label{tab:modern_summary}
\small
\begin{tabular}{lrrrrr}
\toprule
Model & $T_{\rm train}$ & Terminal 8 & Terminal 32 & Early 8 & Loss (pp) \\
\midrule
SwiGLU & 8 & $80.87\pm0.09$ & $75.76\pm0.96$ & $81.22\pm0.05$ & $5.11\pm0.94$ \\
SwiGLU & 32 & $80.43\pm0.05$ & $80.84\pm0.15$ & $79.72\pm0.14$ & $-0.42\pm0.18$ \\
attention & 8 & $79.62\pm0.10$ & $77.20\pm0.34$ & $79.88\pm0.11$ & $2.42\pm0.29$ \\
attention & 32 & $79.20\pm0.13$ & $79.16\pm0.22$ & $78.58\pm0.29$ & $0.04\pm0.17$ \\
\bottomrule
\end{tabular}
\end{table}


\begin{figure}[htbp]
\centering
\includegraphics[width=0.85\linewidth]{fig_modern_horizon_drop.pdf}
\caption{\textbf{Matched training removes the strong horizon-extension loss.}
Terminal $A(8)-A(32)$ for the four recursive conditions. Negative values mean
the endpoint improves with additional steps. Points show individual seeds;
bars and error bars show means and sample standard deviations ($n=3$).
This is the difference summary of the curves in \Cref{fig:modern_panel}.}
\label{fig:modern_drop}
\end{figure}

\section{Tiny ImageNet: Readouts and Representation Probes}
\label{app:modern_readouts}

\paragraph{Frozen readout settings.}
Early-window uses $K=8$, adaptive stopping uses norm-ratio threshold 1.5,
and \grace{} uses $(\tau,\alpha)=(1.2,10)$. These settings are not tuned for
attention. Uniform averaging and \grace{} use the available trajectory through
each evaluation horizon. The shared classification head is applied to the
aggregated state, without retraining. The table below retains all attention
readouts, including underperformance of the norm-weighted rule in the matched
32-step condition. The modern experiments support an endpoint-validity
diagnosis rather than universal dominance of \grace{}.

\begin{table}[htbp]
\centering
\caption{\textbf{All frozen attention readouts.} Validation accuracy (\%, mean $\pm$ sample standard deviation, three seeds).}
\label{tab:attention_readouts}
\small
\begin{tabular}{rrrrrrr}
\toprule
$T_{\rm train}$ & $T_{\rm test}$ & Terminal & Early 8 & Uniform & Adaptive & \grace{} \\
\midrule
8 & 8 & $79.62\pm0.08$ & $79.88\pm0.11$ & $79.88\pm0.11$ & $79.40\pm0.20$ & $77.16\pm0.36$ \\
8 & 16 & $78.85\pm0.20$ & $79.88\pm0.11$ & $79.41\pm0.07$ & $79.40\pm0.20$ & $77.16\pm0.36$ \\
8 & 24 & $78.03\pm0.20$ & $79.88\pm0.11$ & $78.92\pm0.21$ & $79.40\pm0.20$ & $77.16\pm0.36$ \\
8 & 32 & $77.20\pm0.33$ & $79.88\pm0.11$ & $78.49\pm0.19$ & $79.40\pm0.20$ & $77.16\pm0.36$ \\
\midrule
32 & 8 & $79.19\pm0.11$ & $78.58\pm0.29$ & $78.58\pm0.29$ & $76.81\pm0.33$ & $57.47\pm3.94$ \\
32 & 16 & $79.36\pm0.19$ & $78.58\pm0.29$ & $79.40\pm0.16$ & $76.81\pm0.33$ & $57.47\pm3.94$ \\
32 & 24 & $79.32\pm0.37$ & $78.58\pm0.29$ & $79.41\pm0.26$ & $76.81\pm0.33$ & $57.47\pm3.94$ \\
32 & 32 & $79.15\pm0.22$ & $78.58\pm0.29$ & $79.41\pm0.24$ & $76.81\pm0.33$ & $57.47\pm3.94$ \\
\bottomrule
\end{tabular}
\end{table}


The horizon headline uses the saved full-trajectory classification results;
the readout table uses the separately classified feature aggregates. Their
attention terminal scores differ by a few examples because BF16 classifier
matrix multiplications use different tensor shapes. Each series is retained
consistently; these small differences do not alter the reported pattern.

\paragraph{Fixed-regularisation representation probes.}
We fit separate RidgeClassifiers with $\alpha_{\rm ridge}=1$ on step-8,
early-eight, and step-32 features, using the same fixed 20,000-example
training subset (subset seed 2026), no feature standardisation, and all
10,000 validation examples. The regularisation is fixed, while each
representation receives its own fitted coefficients. This distinguishes a
loss of linear decodability from failure of the shared trained head alone.

\begin{table}[htbp]
\centering
\caption{\textbf{Fixed-alpha probe accuracy.} Separately fitted probes share ridge $\alpha=1$ and a fixed training subset. Accuracy in percent; loss in pp.}
\label{tab:modern_probes}
\small
\begin{tabular}{lrrrrr}
\toprule
Model & $T_{\rm train}$ & Step 8 & Early 8 & Step 32 & Loss (pp) \\
\midrule
SwiGLU & 8 & $81.31\pm0.14$ & $81.48\pm0.12$ & $79.08\pm0.29$ & $2.23\pm0.17$ \\
SwiGLU & 32 & $81.12\pm0.01$ & $80.75\pm0.06$ & $81.36\pm0.09$ & $-0.24\pm0.08$ \\
Attention & 8 & $79.59\pm0.06$ & $79.79\pm0.07$ & $78.31\pm0.08$ & $1.28\pm0.08$ \\
Attention & 32 & $79.33\pm0.35$ & $79.39\pm0.30$ & $78.98\pm0.30$ & $0.35\pm0.10$ \\
\bottomrule
\end{tabular}
\end{table}


For eight-step training, both architectures lose accuracy under both the
shared head and refitted probes. Matching the training horizon removes the
large probe loss: the SwiGLU probe improves by $0.24\pm0.08$~pp and the
attention probe loses only $0.35\pm0.10$~pp. This is evidence about the tested
linear decoders, not a proof that all task information is irrecoverably lost.

\section{Tiny ImageNet: Robustness and Preliminary Controls}
\label{app:modern_robustness}

\paragraph{Classwise robustness.}
The classwise loss is the mean step-8 accuracy minus mean step-32 accuracy
across seeds, evaluated separately in each of the 200 classes. All classes
are included. SwiGLU has a median loss of 2.67~pp (IQR 0.00--8.67), with
144/200 classes losing accuracy and 124/200 losing at least 2~pp. Attention
has a median loss of 2.00~pp (IQR 0.67--4.00), with 153/200 classes losing
accuracy and 105/200 losing at least 2~pp. The macro-mean losses are 5.11
and 2.42~pp, respectively. The degradation is therefore distributed across
many classes.

\paragraph{Checkpoint sensitivity.}
We evaluate existing saved states without changing the headline selection
rule. Every audited checkpoint has a positive step-8-minus-step-32 loss,
but its magnitude depends on checkpoint age. SwiGLU's selected checkpoints
(epochs 3--4) lose 4.27--6.13~pp, while final-epoch checkpoints lose only
0.59--0.68~pp. Continued eight-step training therefore substantially reduces
the loss even without matched 32-step training. For attention, the top three
saved improving checkpoints per seed lose 2.12--3.49~pp; the selected
checkpoint has the smallest loss among those saved alternatives in each seed.
Intermediate SwiGLU checkpoints and final-epoch attention states were not
retained, limiting the comparisons possible from saved artifacts.

\begin{table}[htbp]
\centering
\caption{\textbf{Checkpoint sensitivity of eight-step-trained models.} Accuracy in percent; loss in pp. Only retained checkpoints are compared.}
\label{tab:checkpoint_sensitivity}
\small
\begin{tabular}{lr l rrrr}
\toprule
Model & Seed & Checkpoint & Epoch & Step 8 & Step 32 & Loss \\
\midrule
SwiGLU & 101 & Selected & 3 & 80.98 & 76.05 & 4.93 \\
SwiGLU & 101 & Final epoch & 30 & 80.71 & 80.12 & 0.59 \\
SwiGLU & 102 & Selected & 4 & 80.82 & 74.69 & 6.13 \\
SwiGLU & 102 & Final epoch & 30 & 80.62 & 79.96 & 0.66 \\
SwiGLU & 103 & Selected & 3 & 80.82 & 76.55 & 4.27 \\
SwiGLU & 103 & Final epoch & 30 & 80.67 & 79.99 & 0.68 \\
Attention & 101 & Selected & 5 & 79.53 & 76.83 & 2.70 \\
Attention & 101 & Saved rank 2 & 4 & 78.91 & 76.08 & 2.83 \\
Attention & 101 & Saved rank 3 & 3 & 78.52 & 75.03 & 3.49 \\
Attention & 102 & Selected & 7 & 79.61 & 77.49 & 2.12 \\
Attention & 102 & Saved rank 2 & 5 & 79.31 & 76.76 & 2.55 \\
Attention & 102 & Saved rank 3 & 4 & 79.27 & 76.28 & 2.99 \\
Attention & 103 & Selected & 5 & 79.72 & 77.27 & 2.45 \\
Attention & 103 & Saved rank 2 & 4 & 79.30 & 76.63 & 2.67 \\
Attention & 103 & Saved rank 3 & 3 & 78.38 & 75.06 & 3.32 \\
\bottomrule
\end{tabular}
\end{table}


\paragraph{Seed-0 gain-sweep pilots.}
Before the paired panel, two SwiGLU pilots used gain values
$\{0.75,1.0,1.1,1.2,1.3,1.4,1.5,1.6\}$. The first trained for eight steps
and evaluated for 32. Its early-minus-terminal gap was already 4.93~pp at
nominal gain, so it failed the predeclared requirement for a competitive
nominal endpoint before gain stress. The aligned 32-step pilot restored the
nominal endpoint but had a maximum early-window advantage of only 0.27~pp,
failing the criterion of at least 1~pp at three consecutive gains. Each pilot
was stopped after seed 0; neither is a three-seed confirmation. These outcomes
motivated the subsequently frozen paired horizon comparison with fresh seeds.
The direct fixed-alpha probe on pooled ConvNeXt features scored 80.52\%.

\begin{table}[htbp]
\centering
\caption{\textbf{Seed-0 SwiGLU pilot readouts at $T_{\rm test}=32$.} Accuracy in percent. Fixed denotes the training-derived global \grace{} profile. Neither pilot continued to further seeds.}
\label{tab:pilot_readouts}
\small
\begin{tabular}{rrrrrrrr}
\toprule
$T_{\rm train}$ & Gain & Terminal & Early 8 & Uniform & Adaptive & \grace{} & Fixed \\
\midrule
8 & 0.75 & 77.72 & 80.76 & 79.71 & 80.34 & 80.70 & 80.36 \\
8 & 1.00 & 75.92 & 80.85 & 78.84 & 80.38 & 80.70 & 80.39 \\
8 & 1.10 & 75.09 & 80.89 & 78.52 & 80.35 & 80.69 & 80.56 \\
8 & 1.20 & 74.47 & 80.78 & 78.31 & 80.36 & 80.72 & 80.59 \\
8 & 1.30 & 73.77 & 80.82 & 77.92 & 80.32 & 80.68 & 80.68 \\
8 & 1.40 & 73.08 & 80.87 & 77.65 & 80.28 & 80.71 & 80.78 \\
8 & 1.50 & 72.44 & 80.83 & 77.41 & 80.37 & 80.69 & 80.80 \\
8 & 1.60 & 71.86 & 80.82 & 77.13 & 80.29 & 80.76 & 80.81 \\
\midrule
32 & 0.75 & 80.77 & 79.25 & 80.63 & 80.77 & 80.50 & 80.06 \\
32 & 1.00 & 80.75 & 79.72 & 80.85 & 80.71 & 80.64 & 80.33 \\
32 & 1.10 & 80.62 & 79.77 & 80.78 & 80.64 & 80.52 & 80.33 \\
32 & 1.20 & 80.54 & 79.84 & 80.77 & 80.59 & 80.54 & 80.39 \\
32 & 1.30 & 80.33 & 79.95 & 80.71 & 80.57 & 80.59 & 80.40 \\
32 & 1.40 & 80.15 & 80.07 & 80.80 & 80.43 & 80.63 & 80.51 \\
32 & 1.50 & 80.15 & 80.15 & 80.77 & 80.57 & 80.67 & 80.70 \\
32 & 1.60 & 79.93 & 80.20 & 80.79 & 80.49 & 80.66 & 80.68 \\
\bottomrule
\end{tabular}
\end{table}


The pilots also compare per-sample \grace{} with a fixed temporal profile,
formed by averaging normalised training-sample \grace{} weights at gain 1.3
and then freezing that profile across gains. They show no consistent advantage
for per-sample weighting. At gain 1.6 in the aligned pilot, uniform averaging,
\grace{}, and the fixed profile score 80.79\%, 80.66\%, and 80.68\%,
respectively. A saved-trajectory horizon audit of the eight-step pilot at gain
1 gives terminal accuracy 80.82\% at step 8 and 75.92\% at step 32, consistent
with the later paired panel. No new canonical-DEQ experiment was completed;
the DEQ-style evidence in \Cref{app:deq} remains the original limited check.

\section{Tiny ImageNet: Model Scale, Runtime, and Memory}
\label{app:modern_cost}

\begin{table}[htbp]
\centering
\caption{\textbf{Scale of the modern recursive systems.} The executed frozen
backbone has 27,820,128 parameters. One multiply-add counts as two FLOPs.}
\label{tab:modern_scale}
\small
\begin{tabular}{lrr}
\toprule
Quantity & SwiGLU & Attention \\
\midrule
Input & pooled 768-d & $49\times768$ tokens \\
Recursive width & 512 & 256 \\
Intermediate width & 2,048 & 1,024 \\
Attention heads & --- & 8 \\
Trainable head parameters & 3,648,712 & 1,301,960 \\
Executed total parameters & 31,468,840 & 29,122,088 \\
Dominant recurrent FLOPs/step & 6,291,456 & 107,517,600 \\
\bottomrule
\end{tabular}
\end{table}

For state width $d$, intermediate width $m$, $n$ tokens, and $h$ attention
heads, recurrent FLOPs are counted as $6dm$ for SwiGLU and
$8nd^2+4n^2d+5hn^2+6ndm$ for attention ($n=50$ including the class token).
These counts exclude the frozen backbone, one-time input projection, final
classifier, LayerNorm, and lower-order elementwise operations.

\paragraph{Readout benchmark.}
The measurements in \Cref{tab:cost_main} use the first seed-0 SwiGLU pilot,
independently of its scientific continuation outcome. Settings are BF16,
batch 512, $d=512$, $T=32$, gain 1.3, and early window 8, on an RTX 5080
Laptop GPU. CUDA events measure elapsed device time after 10 warm-ups,
with the median of 25 timed repetitions reported. Full recurrent-head
inference starts with cached backbone features and includes recursion and
classification; it excludes image loading and backbone execution. Additional
peak allocated memory is measured relative to the warmed-up baseline for
each repetition and summarised by its median. It excludes baseline model,
input, and cached-tensor allocations and is not CUDA reserved or total process
memory. Readout-only measurements start with an already allocated trajectory.

\begin{table}[htbp]
\centering
\caption{\textbf{Full recurrent-head and readout-only benchmark.} Latencies in ms and additional peak allocations in MiB; the precomputed trajectory is part of the baseline in readout-only measurements.}
\label{tab:cost_detail}
\small
\begin{tabular}{lrrrr}
\toprule
Readout & Full ms & Readout ms & Full MiB & Readout MiB \\
\midrule
Terminal & 6.727 & 0.113 & 8.00 & 2.00 \\
Early window & 1.803 & 0.135 & 11.50 & 2.50 \\
Adaptive & 13.840 & 0.208 & 9.01 & 2.52 \\
Uniform streaming & 6.889 & 0.117 & 8.50 & 2.50 \\
\grace{}, stored & 6.948 & 0.317 & 48.56 & 32.06 \\
\grace{}, streaming & 9.681 & 3.653 & 8.51 & 3.01 \\
\bottomrule
\end{tabular}
\end{table}


\begin{table}[htbp]
\centering
\caption{\textbf{Readout complexity per example.} Arithmetic excludes the
recurrent update and classifier. $t\leq T$ is the adaptive stopping step.}
\label{tab:readout_complexity}
\small
\begin{tabular}{lccc}
\toprule
Readout & Arithmetic & State storage & Halts recurrence \\
\midrule
Terminal & $O(d)$ & $O(d)$ & No \\
Early window & $O(Kd)$ & $O(d)$ & At $K$ \\
Adaptive & $O(td)$ & $O(d)$ & Yes \\
Uniform streaming & $O(Td)$ & $O(d)$ & No \\
\grace{}, stored & $O(Td)$ & $O(Td)$ & No \\
\grace{}, streaming & $O(Td)$ & $O(d)$ & No \\
\bottomrule
\end{tabular}
\end{table}

Streaming \grace{} maintains a weighted numerator and denominator while
discarding each state after use. This explains its memory benefit without
implying a speedup: the measured eager loop costs 9.68~ms versus 6.95~ms
for stored-trajectory \grace{}. Adaptive active-set compaction costs 13.84~ms
for this dense batch. These are implementation- and hardware-specific
measurements, not architecture-independent latency guarantees.

\paragraph{Training and evaluation allocations.}
The attention training runs separately record absolute CUDA allocator peaks.
For eight-step training, allocated/reserved peaks are 871.56/1,100~MiB;
for 32-step training they are 3,054.91/3,594~MiB. Evaluation peaks are
219.53/298~MiB in both conditions. The peaks are identical across the three
seeds within each condition. Completed training wall times are
499.42/498.15/505.00~s for eight steps and
1,582.40/1,635.76/2,312.70~s for 32 steps. These absolute allocator peaks
measure a different quantity from the additional-memory benchmark above.

\paragraph{Execution provenance.}
Protocols, configurations, seeds, and checkpoint selection were frozen for
each experiment before its results. The final panel contains 12 completed
scientific runs, plus the two separate seed-0 pilots. Two interrupted
32-step attention attempts were retained and restarted from the same frozen
seed/configuration because optimizer and RNG states were unavailable for an
exact resume; they are not extra scientific replicates. Scientific artifacts
retain SHA-256 configuration and checkpoint records. The robustness analyses
use existing checkpoints, with no post-result changes to architectures,
gains, readout thresholds, probe regularisation, or the headline selection rule.
